% =====================================================================
% New subsection for Section 2 ("The credible-subpopulation LATT").
% Suggested placement: immediately after \subsection{Honest inference on
% the selected set} (i.e. after Lemma~\ref{lem:calib} and the
% \Delta^{RM} caution, just before \section{Simulation study}), as
% \subsection{When the trade pays: a dominance condition}\label{sec:dominance}.
% Insert with % =====================================================================
% New subsection for Section 2 ("The credible-subpopulation LATT").
% Suggested placement: immediately after \subsection{Honest inference on
% the selected set} (i.e. after Lemma~\ref{lem:calib} and the
% \Delta^{RM} caution, just before \section{Simulation study}), as
% \subsection{When the trade pays: a dominance condition}\label{sec:dominance}.
% Insert with % =====================================================================
% New subsection for Section 2 ("The credible-subpopulation LATT").
% Suggested placement: immediately after \subsection{Honest inference on
% the selected set} (i.e. after Lemma~\ref{lem:calib} and the
% \Delta^{RM} caution, just before \section{Simulation study}), as
% \subsection{When the trade pays: a dominance condition}\label{sec:dominance}.
% Insert with % =====================================================================
% New subsection for Section 2 ("The credible-subpopulation LATT").
% Suggested placement: immediately after \subsection{Honest inference on
% the selected set} (i.e. after Lemma~\ref{lem:calib} and the
% \Delta^{RM} caution, just before \section{Simulation study}), as
% \subsection{When the trade pays: a dominance condition}\label{sec:dominance}.
% Insert with \input{dominance_section} or paste inline.
% Requires the existing theorem environments (proposition, and add
% \newtheorem{corollary}{Corollary} to the preamble if not present).
% =====================================================================

\subsection{When the trade pays: a dominance condition}\label{sec:dominance}

The scope condition documented in the simulations below has a decision-theoretic
counterpart that can be stated in closed form. Changing the estimand is a trade, and
a trade is worth making only when what it buys exceeds what it costs. We make the
ledger explicit in the leading two-cohort case and show that its sign is governed by
the size of the violation relative to the noise and the treatment-effect
heterogeneity, and---crucially---not by the informativeness of the pre-trends, which
governs only how much of a fixed-sign advantage is realized.

Consider two equally weighted cohorts sharing a common sampling variance $\sigma^2$,
observed as one pre- and one post-treatment coefficient each, with independent pre-
and post-period sampling errors. One cohort is clean, $\delta_1=0$; the other is
confounded, with a post-treatment violation $V>0$ and a pre-treatment violation
$\phi V$, where $\phi\ge 0$ is the informativeness of the pre-trend. Their treatment
effects are $\tau_1$ and $\tau_2$, with $\gamma:=\tau_2-\tau_1$ the effect
heterogeneity, so the target is the average $\bar\tau=(\tau_1+\tau_2)/2$. The screen
retains a cohort when $\lvert\hat\beta_{g,\mathrm{pre}}\rvert\le c$; we take $c$ large
enough that the clean cohort is retained with probability approaching one, and write
\begin{equation}\label{eq:p2}
p_2(\phi)\;=\;\Pr\!\big(\lvert\hat\beta_{2,\mathrm{pre}}\rvert\le c\big)
\;=\;\Phi\!\Big(\tfrac{c-\phi V}{\sigma}\Big)-\Phi\!\Big(\tfrac{-c-\phi V}{\sigma}\Big)
\end{equation}
for the probability that the confounded cohort survives the screen, which decreases
from near one, when the confound is invisible in the pre-period, to zero as the
pre-trend announces it. The ATT estimator averages the post-treatment coefficients of
both cohorts; the LATT estimator averages over the retained set, falling back to both
cohorts when neither survives. We compare them by mean squared error for the common
target $\bar\tau$, under which the two estimands coincide and the comparison is a fair
one between procedures aimed at the same number.

\begin{proposition}[When the trade pays]\label{prop:dominance}
Under the two-cohort model above, the mean-squared-error advantage of the LATT
estimator over the ATT estimator is
\begin{equation}\label{eq:delta}
\Delta(\phi)\;=\;\mathrm{MSE}_{\mathrm{ATT}}-\mathrm{MSE}_{\mathrm{LATT}}
\;=\;\frac{1-p_2(\phi)}{4}\,\big[\,V^2-\gamma^2-2\sigma^2\,\big].
\end{equation}
Consequently the LATT estimator dominates the ATT estimator in mean squared error if
and only if
\begin{equation}\label{eq:threshold}
V^2\;>\;\gamma^2+2\sigma^2 .
\end{equation}
The sign of the advantage is fixed by \eqref{eq:threshold} and is independent of the
informativeness $\phi$; informativeness enters only through the factor $1-p_2(\phi)$,
which is nondecreasing in $\phi$ and scales the magnitude of a fixed-sign advantage
from zero, when the pre-trend is uninformative, to its ceiling $\tfrac14[V^2-\gamma^2-2\sigma^2]$
as $\phi\to\infty$.
\end{proposition}

\begin{proof}
Write $\tau_1=\bar\tau-\gamma/2$ and $\tau_2=\bar\tau+\gamma/2$. The post-treatment
coefficients are $\hat\beta_{1,\mathrm{post}}\sim\mathcal N(\bar\tau-\gamma/2,\sigma^2)$
and $\hat\beta_{2,\mathrm{post}}\sim\mathcal N(\bar\tau+\gamma/2+V,\sigma^2)$,
independent of each other and, by the independence of pre- and post-period errors, of
the retention indicators. The ATT estimator
$\tfrac12(\hat\beta_{1,\mathrm{post}}+\hat\beta_{2,\mathrm{post}})$ has mean
$\bar\tau+V/2$ and variance $\sigma^2/2$, so
\[
\mathrm{MSE}_{\mathrm{ATT}}=\frac{V^2}{4}+\frac{\sigma^2}{2}.
\]
With the clean cohort retained with probability one, the LATT estimator takes two
values. With probability $p_2$ the confounded cohort is also retained and the
estimator is the two-cohort average, centered at $\bar\tau+V/2$ with variance
$\sigma^2/2$; with probability $1-p_2$ only the clean cohort is retained and the
estimator is $\hat\beta_{1,\mathrm{post}}$, centered at $\bar\tau-\gamma/2$ with
variance $\sigma^2$. Because retention is independent of the post-period errors, the
mean squared error is the retention-probability-weighted average of the two branch
mean squared errors, each of which is its own squared bias plus variance,
\[
\mathrm{MSE}_{\mathrm{LATT}}
=p_2\!\left(\frac{V^2}{4}+\frac{\sigma^2}{2}\right)
+(1-p_2)\!\left(\frac{\gamma^2}{4}+\sigma^2\right).
\]
Subtracting,
\[
\Delta
=(1-p_2)\!\left(\frac{V^2}{4}+\frac{\sigma^2}{2}\right)
-(1-p_2)\!\left(\frac{\gamma^2}{4}+\sigma^2\right)
=\frac{1-p_2}{4}\big[V^2-\gamma^2-2\sigma^2\big],
\]
which is \eqref{eq:delta}. Since $1-p_2(\phi)\ge 0$, the sign of $\Delta$ is that of
the bracket whenever $p_2<1$, giving \eqref{eq:threshold}; and $p_2(\phi)$ in
\eqref{eq:p2} is nonincreasing in $\phi$ for $\phi V\ge 0$, so $1-p_2$ is
nondecreasing, establishing the last claim.
\end{proof}

The bracket in \eqref{eq:delta} is the ledger, and its three terms are the three
forces at work. The $+V^2$ term is the bias the screen removes by dropping the
confounded cohort, the reason to trade. The $-2\sigma^2$ term is the variance the
trade costs, since dropping a cohort halves the effective sample. The $-\gamma^2$ term
is the price of changing the estimand: the retained subpopulation's effect differs
from the ATT by the heterogeneity $\gamma$, so a researcher who reports the LATT while
targeting the ATT inherits that difference as a bias. This last term is the
``different estimand'' objection made quantitative, and it vanishes exactly when
treatment effects are homogeneous, $\gamma=0$, which is the normalization adopted in
the simulations. The dominance condition thus reveals that normalization to be the
most favorable case for the method, and names precisely what departing from it costs.

Two readings of \eqref{eq:threshold} bear on the paper's message. First, informative
pre-trends are necessary but not sufficient for the trade to pay. If
$V^2<\gamma^2+2\sigma^2$---a violation small relative to the noise and the
heterogeneity---then no degree of informativeness rescues the LATT: as $\phi$ grows the
screen correctly detects and discards the confounded cohort, yet the factor $1-p_2$
merely drives an already-negative bracket further from zero, and the LATT's mean
squared error is strictly larger. In that regime the efficient decision is to retain
the biased cohort, because the bias it carries contributes less to risk than the
variance its removal forgoes. Informativeness governs how fully the advantage is
realized, through $1-p_2$; it does not create one. Second, when the violation is large
enough that the bracket is positive, the advantage rises monotonically with
informativeness from nothing, at $\phi=0$ where the two curves coincide, to a ceiling
set by the bracket---which is the vertical gap between the two point-estimate bias
curves traced in Figure~\ref{fig:scope} read as a function of $\phi$.

The same accounting extends to $K$ cohorts, at the cost of a closed form. Partition
the cohorts into a clean block, retained with probability approaching one, and
confounded cohorts each admitted with a probability $\pi_g$ that the screen drives
toward zero as its pre-trend becomes informative. To first order in the admission
probabilities the LATT's risk is a mixture, dominated by the event that no confounded
cohort survives and, at order $\pi_g$, by the single-admission events, and the sign of
the leading advantage is again governed by a net-value functional in which each
retained-versus-dropped comparison contributes a violation term against a variance
term and a heterogeneity term. The two-cohort ledger is therefore the mechanism in
miniature, and the qualitative conclusion---that the sign is a violation-to-noise
question and informativeness a magnitude question---does not depend on the cohort
count.

Three caveats keep the result honest. The advantage \eqref{eq:delta} is an upper bound
on what the method delivers, because it omits two channels that both work against the
LATT: occasional false rejection of the clean cohort, which enters at order
$1-p_1$ when the threshold is not comfortably large, and the additional variance that
data-driven selection induces and that a standard error conditional on the realized
set does not capture, the understatement documented in Section~\ref{sec:calib}. Both
shrink the realized advantage without changing its sign or its shape. Finally,
\eqref{eq:delta} is a comparison of point estimators under squared-error loss, not of
the honest intervals that are the paper's recommended output. The interval analogue
replaces the reason-to-trade $V$ with the residual violation net of the calibration
buffer of Lemma~\ref{lem:calib}, and the variance cost $\sigma^2$ with a sampling
slack inflated by the selection quantile, so that the level-bound half-width comparison
turns on the same three forces with $V$ discounted and $\sigma$ inflated. The
mean-squared-error threshold \eqref{eq:threshold} correctly signs that comparison and
is the transparent summary; the interval version sharpens the constants.
 or paste inline.
% Requires the existing theorem environments (proposition, and add
% \newtheorem{corollary}{Corollary} to the preamble if not present).
% =====================================================================

\subsection{When the trade pays: a dominance condition}\label{sec:dominance}

The scope condition documented in the simulations below has a decision-theoretic
counterpart that can be stated in closed form. Changing the estimand is a trade, and
a trade is worth making only when what it buys exceeds what it costs. We make the
ledger explicit in the leading two-cohort case and show that its sign is governed by
the size of the violation relative to the noise and the treatment-effect
heterogeneity, and---crucially---not by the informativeness of the pre-trends, which
governs only how much of a fixed-sign advantage is realized.

Consider two equally weighted cohorts sharing a common sampling variance $\sigma^2$,
observed as one pre- and one post-treatment coefficient each, with independent pre-
and post-period sampling errors. One cohort is clean, $\delta_1=0$; the other is
confounded, with a post-treatment violation $V>0$ and a pre-treatment violation
$\phi V$, where $\phi\ge 0$ is the informativeness of the pre-trend. Their treatment
effects are $\tau_1$ and $\tau_2$, with $\gamma:=\tau_2-\tau_1$ the effect
heterogeneity, so the target is the average $\bar\tau=(\tau_1+\tau_2)/2$. The screen
retains a cohort when $\lvert\hat\beta_{g,\mathrm{pre}}\rvert\le c$; we take $c$ large
enough that the clean cohort is retained with probability approaching one, and write
\begin{equation}\label{eq:p2}
p_2(\phi)\;=\;\Pr\!\big(\lvert\hat\beta_{2,\mathrm{pre}}\rvert\le c\big)
\;=\;\Phi\!\Big(\tfrac{c-\phi V}{\sigma}\Big)-\Phi\!\Big(\tfrac{-c-\phi V}{\sigma}\Big)
\end{equation}
for the probability that the confounded cohort survives the screen, which decreases
from near one, when the confound is invisible in the pre-period, to zero as the
pre-trend announces it. The ATT estimator averages the post-treatment coefficients of
both cohorts; the LATT estimator averages over the retained set, falling back to both
cohorts when neither survives. We compare them by mean squared error for the common
target $\bar\tau$, under which the two estimands coincide and the comparison is a fair
one between procedures aimed at the same number.

\begin{proposition}[When the trade pays]\label{prop:dominance}
Under the two-cohort model above, the mean-squared-error advantage of the LATT
estimator over the ATT estimator is
\begin{equation}\label{eq:delta}
\Delta(\phi)\;=\;\mathrm{MSE}_{\mathrm{ATT}}-\mathrm{MSE}_{\mathrm{LATT}}
\;=\;\frac{1-p_2(\phi)}{4}\,\big[\,V^2-\gamma^2-2\sigma^2\,\big].
\end{equation}
Consequently the LATT estimator dominates the ATT estimator in mean squared error if
and only if
\begin{equation}\label{eq:threshold}
V^2\;>\;\gamma^2+2\sigma^2 .
\end{equation}
The sign of the advantage is fixed by \eqref{eq:threshold} and is independent of the
informativeness $\phi$; informativeness enters only through the factor $1-p_2(\phi)$,
which is nondecreasing in $\phi$ and scales the magnitude of a fixed-sign advantage
from zero, when the pre-trend is uninformative, to its ceiling $\tfrac14[V^2-\gamma^2-2\sigma^2]$
as $\phi\to\infty$.
\end{proposition}

\begin{proof}
Write $\tau_1=\bar\tau-\gamma/2$ and $\tau_2=\bar\tau+\gamma/2$. The post-treatment
coefficients are $\hat\beta_{1,\mathrm{post}}\sim\mathcal N(\bar\tau-\gamma/2,\sigma^2)$
and $\hat\beta_{2,\mathrm{post}}\sim\mathcal N(\bar\tau+\gamma/2+V,\sigma^2)$,
independent of each other and, by the independence of pre- and post-period errors, of
the retention indicators. The ATT estimator
$\tfrac12(\hat\beta_{1,\mathrm{post}}+\hat\beta_{2,\mathrm{post}})$ has mean
$\bar\tau+V/2$ and variance $\sigma^2/2$, so
\[
\mathrm{MSE}_{\mathrm{ATT}}=\frac{V^2}{4}+\frac{\sigma^2}{2}.
\]
With the clean cohort retained with probability one, the LATT estimator takes two
values. With probability $p_2$ the confounded cohort is also retained and the
estimator is the two-cohort average, centered at $\bar\tau+V/2$ with variance
$\sigma^2/2$; with probability $1-p_2$ only the clean cohort is retained and the
estimator is $\hat\beta_{1,\mathrm{post}}$, centered at $\bar\tau-\gamma/2$ with
variance $\sigma^2$. Because retention is independent of the post-period errors, the
mean squared error is the retention-probability-weighted average of the two branch
mean squared errors, each of which is its own squared bias plus variance,
\[
\mathrm{MSE}_{\mathrm{LATT}}
=p_2\!\left(\frac{V^2}{4}+\frac{\sigma^2}{2}\right)
+(1-p_2)\!\left(\frac{\gamma^2}{4}+\sigma^2\right).
\]
Subtracting,
\[
\Delta
=(1-p_2)\!\left(\frac{V^2}{4}+\frac{\sigma^2}{2}\right)
-(1-p_2)\!\left(\frac{\gamma^2}{4}+\sigma^2\right)
=\frac{1-p_2}{4}\big[V^2-\gamma^2-2\sigma^2\big],
\]
which is \eqref{eq:delta}. Since $1-p_2(\phi)\ge 0$, the sign of $\Delta$ is that of
the bracket whenever $p_2<1$, giving \eqref{eq:threshold}; and $p_2(\phi)$ in
\eqref{eq:p2} is nonincreasing in $\phi$ for $\phi V\ge 0$, so $1-p_2$ is
nondecreasing, establishing the last claim.
\end{proof}

The bracket in \eqref{eq:delta} is the ledger, and its three terms are the three
forces at work. The $+V^2$ term is the bias the screen removes by dropping the
confounded cohort, the reason to trade. The $-2\sigma^2$ term is the variance the
trade costs, since dropping a cohort halves the effective sample. The $-\gamma^2$ term
is the price of changing the estimand: the retained subpopulation's effect differs
from the ATT by the heterogeneity $\gamma$, so a researcher who reports the LATT while
targeting the ATT inherits that difference as a bias. This last term is the
``different estimand'' objection made quantitative, and it vanishes exactly when
treatment effects are homogeneous, $\gamma=0$, which is the normalization adopted in
the simulations. The dominance condition thus reveals that normalization to be the
most favorable case for the method, and names precisely what departing from it costs.

Two readings of \eqref{eq:threshold} bear on the paper's message. First, informative
pre-trends are necessary but not sufficient for the trade to pay. If
$V^2<\gamma^2+2\sigma^2$---a violation small relative to the noise and the
heterogeneity---then no degree of informativeness rescues the LATT: as $\phi$ grows the
screen correctly detects and discards the confounded cohort, yet the factor $1-p_2$
merely drives an already-negative bracket further from zero, and the LATT's mean
squared error is strictly larger. In that regime the efficient decision is to retain
the biased cohort, because the bias it carries contributes less to risk than the
variance its removal forgoes. Informativeness governs how fully the advantage is
realized, through $1-p_2$; it does not create one. Second, when the violation is large
enough that the bracket is positive, the advantage rises monotonically with
informativeness from nothing, at $\phi=0$ where the two curves coincide, to a ceiling
set by the bracket---which is the vertical gap between the two point-estimate bias
curves traced in Figure~\ref{fig:scope} read as a function of $\phi$.

The same accounting extends to $K$ cohorts, at the cost of a closed form. Partition
the cohorts into a clean block, retained with probability approaching one, and
confounded cohorts each admitted with a probability $\pi_g$ that the screen drives
toward zero as its pre-trend becomes informative. To first order in the admission
probabilities the LATT's risk is a mixture, dominated by the event that no confounded
cohort survives and, at order $\pi_g$, by the single-admission events, and the sign of
the leading advantage is again governed by a net-value functional in which each
retained-versus-dropped comparison contributes a violation term against a variance
term and a heterogeneity term. The two-cohort ledger is therefore the mechanism in
miniature, and the qualitative conclusion---that the sign is a violation-to-noise
question and informativeness a magnitude question---does not depend on the cohort
count.

Three caveats keep the result honest. The advantage \eqref{eq:delta} is an upper bound
on what the method delivers, because it omits two channels that both work against the
LATT: occasional false rejection of the clean cohort, which enters at order
$1-p_1$ when the threshold is not comfortably large, and the additional variance that
data-driven selection induces and that a standard error conditional on the realized
set does not capture, the understatement documented in Section~\ref{sec:calib}. Both
shrink the realized advantage without changing its sign or its shape. Finally,
\eqref{eq:delta} is a comparison of point estimators under squared-error loss, not of
the honest intervals that are the paper's recommended output. The interval analogue
replaces the reason-to-trade $V$ with the residual violation net of the calibration
buffer of Lemma~\ref{lem:calib}, and the variance cost $\sigma^2$ with a sampling
slack inflated by the selection quantile, so that the level-bound half-width comparison
turns on the same three forces with $V$ discounted and $\sigma$ inflated. The
mean-squared-error threshold \eqref{eq:threshold} correctly signs that comparison and
is the transparent summary; the interval version sharpens the constants.
 or paste inline.
% Requires the existing theorem environments (proposition, and add
% \newtheorem{corollary}{Corollary} to the preamble if not present).
% =====================================================================

\subsection{When the trade pays: a dominance condition}\label{sec:dominance}

The scope condition documented in the simulations below has a decision-theoretic
counterpart that can be stated in closed form. Changing the estimand is a trade, and
a trade is worth making only when what it buys exceeds what it costs. We make the
ledger explicit in the leading two-cohort case and show that its sign is governed by
the size of the violation relative to the noise and the treatment-effect
heterogeneity, and---crucially---not by the informativeness of the pre-trends, which
governs only how much of a fixed-sign advantage is realized.

Consider two equally weighted cohorts sharing a common sampling variance $\sigma^2$,
observed as one pre- and one post-treatment coefficient each, with independent pre-
and post-period sampling errors. One cohort is clean, $\delta_1=0$; the other is
confounded, with a post-treatment violation $V>0$ and a pre-treatment violation
$\phi V$, where $\phi\ge 0$ is the informativeness of the pre-trend. Their treatment
effects are $\tau_1$ and $\tau_2$, with $\gamma:=\tau_2-\tau_1$ the effect
heterogeneity, so the target is the average $\bar\tau=(\tau_1+\tau_2)/2$. The screen
retains a cohort when $\lvert\hat\beta_{g,\mathrm{pre}}\rvert\le c$; we take $c$ large
enough that the clean cohort is retained with probability approaching one, and write
\begin{equation}\label{eq:p2}
p_2(\phi)\;=\;\Pr\!\big(\lvert\hat\beta_{2,\mathrm{pre}}\rvert\le c\big)
\;=\;\Phi\!\Big(\tfrac{c-\phi V}{\sigma}\Big)-\Phi\!\Big(\tfrac{-c-\phi V}{\sigma}\Big)
\end{equation}
for the probability that the confounded cohort survives the screen, which decreases
from near one, when the confound is invisible in the pre-period, to zero as the
pre-trend announces it. The ATT estimator averages the post-treatment coefficients of
both cohorts; the LATT estimator averages over the retained set, falling back to both
cohorts when neither survives. We compare them by mean squared error for the common
target $\bar\tau$, under which the two estimands coincide and the comparison is a fair
one between procedures aimed at the same number.

\begin{proposition}[When the trade pays]\label{prop:dominance}
Under the two-cohort model above, the mean-squared-error advantage of the LATT
estimator over the ATT estimator is
\begin{equation}\label{eq:delta}
\Delta(\phi)\;=\;\mathrm{MSE}_{\mathrm{ATT}}-\mathrm{MSE}_{\mathrm{LATT}}
\;=\;\frac{1-p_2(\phi)}{4}\,\big[\,V^2-\gamma^2-2\sigma^2\,\big].
\end{equation}
Consequently the LATT estimator dominates the ATT estimator in mean squared error if
and only if
\begin{equation}\label{eq:threshold}
V^2\;>\;\gamma^2+2\sigma^2 .
\end{equation}
The sign of the advantage is fixed by \eqref{eq:threshold} and is independent of the
informativeness $\phi$; informativeness enters only through the factor $1-p_2(\phi)$,
which is nondecreasing in $\phi$ and scales the magnitude of a fixed-sign advantage
from zero, when the pre-trend is uninformative, to its ceiling $\tfrac14[V^2-\gamma^2-2\sigma^2]$
as $\phi\to\infty$.
\end{proposition}

\begin{proof}
Write $\tau_1=\bar\tau-\gamma/2$ and $\tau_2=\bar\tau+\gamma/2$. The post-treatment
coefficients are $\hat\beta_{1,\mathrm{post}}\sim\mathcal N(\bar\tau-\gamma/2,\sigma^2)$
and $\hat\beta_{2,\mathrm{post}}\sim\mathcal N(\bar\tau+\gamma/2+V,\sigma^2)$,
independent of each other and, by the independence of pre- and post-period errors, of
the retention indicators. The ATT estimator
$\tfrac12(\hat\beta_{1,\mathrm{post}}+\hat\beta_{2,\mathrm{post}})$ has mean
$\bar\tau+V/2$ and variance $\sigma^2/2$, so
\[
\mathrm{MSE}_{\mathrm{ATT}}=\frac{V^2}{4}+\frac{\sigma^2}{2}.
\]
With the clean cohort retained with probability one, the LATT estimator takes two
values. With probability $p_2$ the confounded cohort is also retained and the
estimator is the two-cohort average, centered at $\bar\tau+V/2$ with variance
$\sigma^2/2$; with probability $1-p_2$ only the clean cohort is retained and the
estimator is $\hat\beta_{1,\mathrm{post}}$, centered at $\bar\tau-\gamma/2$ with
variance $\sigma^2$. Because retention is independent of the post-period errors, the
mean squared error is the retention-probability-weighted average of the two branch
mean squared errors, each of which is its own squared bias plus variance,
\[
\mathrm{MSE}_{\mathrm{LATT}}
=p_2\!\left(\frac{V^2}{4}+\frac{\sigma^2}{2}\right)
+(1-p_2)\!\left(\frac{\gamma^2}{4}+\sigma^2\right).
\]
Subtracting,
\[
\Delta
=(1-p_2)\!\left(\frac{V^2}{4}+\frac{\sigma^2}{2}\right)
-(1-p_2)\!\left(\frac{\gamma^2}{4}+\sigma^2\right)
=\frac{1-p_2}{4}\big[V^2-\gamma^2-2\sigma^2\big],
\]
which is \eqref{eq:delta}. Since $1-p_2(\phi)\ge 0$, the sign of $\Delta$ is that of
the bracket whenever $p_2<1$, giving \eqref{eq:threshold}; and $p_2(\phi)$ in
\eqref{eq:p2} is nonincreasing in $\phi$ for $\phi V\ge 0$, so $1-p_2$ is
nondecreasing, establishing the last claim.
\end{proof}

The bracket in \eqref{eq:delta} is the ledger, and its three terms are the three
forces at work. The $+V^2$ term is the bias the screen removes by dropping the
confounded cohort, the reason to trade. The $-2\sigma^2$ term is the variance the
trade costs, since dropping a cohort halves the effective sample. The $-\gamma^2$ term
is the price of changing the estimand: the retained subpopulation's effect differs
from the ATT by the heterogeneity $\gamma$, so a researcher who reports the LATT while
targeting the ATT inherits that difference as a bias. This last term is the
``different estimand'' objection made quantitative, and it vanishes exactly when
treatment effects are homogeneous, $\gamma=0$, which is the normalization adopted in
the simulations. The dominance condition thus reveals that normalization to be the
most favorable case for the method, and names precisely what departing from it costs.

Two readings of \eqref{eq:threshold} bear on the paper's message. First, informative
pre-trends are necessary but not sufficient for the trade to pay. If
$V^2<\gamma^2+2\sigma^2$---a violation small relative to the noise and the
heterogeneity---then no degree of informativeness rescues the LATT: as $\phi$ grows the
screen correctly detects and discards the confounded cohort, yet the factor $1-p_2$
merely drives an already-negative bracket further from zero, and the LATT's mean
squared error is strictly larger. In that regime the efficient decision is to retain
the biased cohort, because the bias it carries contributes less to risk than the
variance its removal forgoes. Informativeness governs how fully the advantage is
realized, through $1-p_2$; it does not create one. Second, when the violation is large
enough that the bracket is positive, the advantage rises monotonically with
informativeness from nothing, at $\phi=0$ where the two curves coincide, to a ceiling
set by the bracket---which is the vertical gap between the two point-estimate bias
curves traced in Figure~\ref{fig:scope} read as a function of $\phi$.

The same accounting extends to $K$ cohorts, at the cost of a closed form. Partition
the cohorts into a clean block, retained with probability approaching one, and
confounded cohorts each admitted with a probability $\pi_g$ that the screen drives
toward zero as its pre-trend becomes informative. To first order in the admission
probabilities the LATT's risk is a mixture, dominated by the event that no confounded
cohort survives and, at order $\pi_g$, by the single-admission events, and the sign of
the leading advantage is again governed by a net-value functional in which each
retained-versus-dropped comparison contributes a violation term against a variance
term and a heterogeneity term. The two-cohort ledger is therefore the mechanism in
miniature, and the qualitative conclusion---that the sign is a violation-to-noise
question and informativeness a magnitude question---does not depend on the cohort
count.

Three caveats keep the result honest. The advantage \eqref{eq:delta} is an upper bound
on what the method delivers, because it omits two channels that both work against the
LATT: occasional false rejection of the clean cohort, which enters at order
$1-p_1$ when the threshold is not comfortably large, and the additional variance that
data-driven selection induces and that a standard error conditional on the realized
set does not capture, the understatement documented in Section~\ref{sec:calib}. Both
shrink the realized advantage without changing its sign or its shape. Finally,
\eqref{eq:delta} is a comparison of point estimators under squared-error loss, not of
the honest intervals that are the paper's recommended output. The interval analogue
replaces the reason-to-trade $V$ with the residual violation net of the calibration
buffer of Lemma~\ref{lem:calib}, and the variance cost $\sigma^2$ with a sampling
slack inflated by the selection quantile, so that the level-bound half-width comparison
turns on the same three forces with $V$ discounted and $\sigma$ inflated. The
mean-squared-error threshold \eqref{eq:threshold} correctly signs that comparison and
is the transparent summary; the interval version sharpens the constants.
 or paste inline.
% Requires the existing theorem environments (proposition, and add
% \newtheorem{corollary}{Corollary} to the preamble if not present).
% =====================================================================

\subsection{When the trade pays: a dominance condition}\label{sec:dominance}

The scope condition documented in the simulations below has a decision-theoretic
counterpart that can be stated in closed form. Changing the estimand is a trade, and
a trade is worth making only when what it buys exceeds what it costs. We make the
ledger explicit in the leading two-cohort case and show that its sign is governed by
the size of the violation relative to the noise and the treatment-effect
heterogeneity, and---crucially---not by the informativeness of the pre-trends, which
governs only how much of a fixed-sign advantage is realized.

Consider two equally weighted cohorts sharing a common sampling variance $\sigma^2$,
observed as one pre- and one post-treatment coefficient each, with independent pre-
and post-period sampling errors. One cohort is clean, $\delta_1=0$; the other is
confounded, with a post-treatment violation $V>0$ and a pre-treatment violation
$\phi V$, where $\phi\ge 0$ is the informativeness of the pre-trend. Their treatment
effects are $\tau_1$ and $\tau_2$, with $\gamma:=\tau_2-\tau_1$ the effect
heterogeneity, so the target is the average $\bar\tau=(\tau_1+\tau_2)/2$. The screen
retains a cohort when $\lvert\hat\beta_{g,\mathrm{pre}}\rvert\le c$; we take $c$ large
enough that the clean cohort is retained with probability approaching one, and write
\begin{equation}\label{eq:p2}
p_2(\phi)\;=\;\Pr\!\big(\lvert\hat\beta_{2,\mathrm{pre}}\rvert\le c\big)
\;=\;\Phi\!\Big(\tfrac{c-\phi V}{\sigma}\Big)-\Phi\!\Big(\tfrac{-c-\phi V}{\sigma}\Big)
\end{equation}
for the probability that the confounded cohort survives the screen, which decreases
from near one, when the confound is invisible in the pre-period, to zero as the
pre-trend announces it. The ATT estimator averages the post-treatment coefficients of
both cohorts; the LATT estimator averages over the retained set, falling back to both
cohorts when neither survives. We compare them by mean squared error for the common
target $\bar\tau$, under which the two estimands coincide and the comparison is a fair
one between procedures aimed at the same number.

\begin{proposition}[When the trade pays]\label{prop:dominance}
Under the two-cohort model above, the mean-squared-error advantage of the LATT
estimator over the ATT estimator is
\begin{equation}\label{eq:delta}
\Delta(\phi)\;=\;\mathrm{MSE}_{\mathrm{ATT}}-\mathrm{MSE}_{\mathrm{LATT}}
\;=\;\frac{1-p_2(\phi)}{4}\,\big[\,V^2-\gamma^2-2\sigma^2\,\big].
\end{equation}
Consequently the LATT estimator dominates the ATT estimator in mean squared error if
and only if
\begin{equation}\label{eq:threshold}
V^2\;>\;\gamma^2+2\sigma^2 .
\end{equation}
The sign of the advantage is fixed by \eqref{eq:threshold} and is independent of the
informativeness $\phi$; informativeness enters only through the factor $1-p_2(\phi)$,
which is nondecreasing in $\phi$ and scales the magnitude of a fixed-sign advantage
from zero, when the pre-trend is uninformative, to its ceiling $\tfrac14[V^2-\gamma^2-2\sigma^2]$
as $\phi\to\infty$.
\end{proposition}

\begin{proof}
Write $\tau_1=\bar\tau-\gamma/2$ and $\tau_2=\bar\tau+\gamma/2$. The post-treatment
coefficients are $\hat\beta_{1,\mathrm{post}}\sim\mathcal N(\bar\tau-\gamma/2,\sigma^2)$
and $\hat\beta_{2,\mathrm{post}}\sim\mathcal N(\bar\tau+\gamma/2+V,\sigma^2)$,
independent of each other and, by the independence of pre- and post-period errors, of
the retention indicators. The ATT estimator
$\tfrac12(\hat\beta_{1,\mathrm{post}}+\hat\beta_{2,\mathrm{post}})$ has mean
$\bar\tau+V/2$ and variance $\sigma^2/2$, so
\[
\mathrm{MSE}_{\mathrm{ATT}}=\frac{V^2}{4}+\frac{\sigma^2}{2}.
\]
With the clean cohort retained with probability one, the LATT estimator takes two
values. With probability $p_2$ the confounded cohort is also retained and the
estimator is the two-cohort average, centered at $\bar\tau+V/2$ with variance
$\sigma^2/2$; with probability $1-p_2$ only the clean cohort is retained and the
estimator is $\hat\beta_{1,\mathrm{post}}$, centered at $\bar\tau-\gamma/2$ with
variance $\sigma^2$. Because retention is independent of the post-period errors, the
mean squared error is the retention-probability-weighted average of the two branch
mean squared errors, each of which is its own squared bias plus variance,
\[
\mathrm{MSE}_{\mathrm{LATT}}
=p_2\!\left(\frac{V^2}{4}+\frac{\sigma^2}{2}\right)
+(1-p_2)\!\left(\frac{\gamma^2}{4}+\sigma^2\right).
\]
Subtracting,
\[
\Delta
=(1-p_2)\!\left(\frac{V^2}{4}+\frac{\sigma^2}{2}\right)
-(1-p_2)\!\left(\frac{\gamma^2}{4}+\sigma^2\right)
=\frac{1-p_2}{4}\big[V^2-\gamma^2-2\sigma^2\big],
\]
which is \eqref{eq:delta}. Since $1-p_2(\phi)\ge 0$, the sign of $\Delta$ is that of
the bracket whenever $p_2<1$, giving \eqref{eq:threshold}; and $p_2(\phi)$ in
\eqref{eq:p2} is nonincreasing in $\phi$ for $\phi V\ge 0$, so $1-p_2$ is
nondecreasing, establishing the last claim.
\end{proof}

The bracket in \eqref{eq:delta} is the ledger, and its three terms are the three
forces at work. The $+V^2$ term is the bias the screen removes by dropping the
confounded cohort, the reason to trade. The $-2\sigma^2$ term is the variance the
trade costs, since dropping a cohort halves the effective sample. The $-\gamma^2$ term
is the price of changing the estimand: the retained subpopulation's effect differs
from the ATT by the heterogeneity $\gamma$, so a researcher who reports the LATT while
targeting the ATT inherits that difference as a bias. This last term is the
``different estimand'' objection made quantitative, and it vanishes exactly when
treatment effects are homogeneous, $\gamma=0$, which is the normalization adopted in
the simulations. The dominance condition thus reveals that normalization to be the
most favorable case for the method, and names precisely what departing from it costs.

Two readings of \eqref{eq:threshold} bear on the paper's message. First, informative
pre-trends are necessary but not sufficient for the trade to pay. If
$V^2<\gamma^2+2\sigma^2$---a violation small relative to the noise and the
heterogeneity---then no degree of informativeness rescues the LATT: as $\phi$ grows the
screen correctly detects and discards the confounded cohort, yet the factor $1-p_2$
merely drives an already-negative bracket further from zero, and the LATT's mean
squared error is strictly larger. In that regime the efficient decision is to retain
the biased cohort, because the bias it carries contributes less to risk than the
variance its removal forgoes. Informativeness governs how fully the advantage is
realized, through $1-p_2$; it does not create one. Second, when the violation is large
enough that the bracket is positive, the advantage rises monotonically with
informativeness from nothing, at $\phi=0$ where the two curves coincide, to a ceiling
set by the bracket---which is the vertical gap between the two point-estimate bias
curves traced in Figure~\ref{fig:scope} read as a function of $\phi$.

The same accounting extends to $K$ cohorts, at the cost of a closed form. Partition
the cohorts into a clean block, retained with probability approaching one, and
confounded cohorts each admitted with a probability $\pi_g$ that the screen drives
toward zero as its pre-trend becomes informative. To first order in the admission
probabilities the LATT's risk is a mixture, dominated by the event that no confounded
cohort survives and, at order $\pi_g$, by the single-admission events, and the sign of
the leading advantage is again governed by a net-value functional in which each
retained-versus-dropped comparison contributes a violation term against a variance
term and a heterogeneity term. The two-cohort ledger is therefore the mechanism in
miniature, and the qualitative conclusion---that the sign is a violation-to-noise
question and informativeness a magnitude question---does not depend on the cohort
count.

Three caveats keep the result honest. The advantage \eqref{eq:delta} is an upper bound
on what the method delivers, because it omits two channels that both work against the
LATT: occasional false rejection of the clean cohort, which enters at order
$1-p_1$ when the threshold is not comfortably large, and the additional variance that
data-driven selection induces and that a standard error conditional on the realized
set does not capture, the understatement documented in Section~\ref{sec:calib}. Both
shrink the realized advantage without changing its sign or its shape. Finally,
\eqref{eq:delta} is a comparison of point estimators under squared-error loss, not of
the honest intervals that are the paper's recommended output. The interval analogue
replaces the reason-to-trade $V$ with the residual violation net of the calibration
buffer of Lemma~\ref{lem:calib}, and the variance cost $\sigma^2$ with a sampling
slack inflated by the selection quantile, so that the level-bound half-width comparison
turns on the same three forces with $V$ discounted and $\sigma$ inflated. The
mean-squared-error threshold \eqref{eq:threshold} correctly signs that comparison and
is the transparent summary; the interval version sharpens the constants.
